\section{Set-theoretic models and consistency bounds}\label{model:section}

The analytic implications concern Banach spaces inside a fixed model of
set theory.  We now verify that their hypotheses give the two required
relative-consistency statements.  Choice is retained throughout.

Real-valued measurability and its strict additivity convention are as defined
in the introduction; this agrees with \cite[Definition~1.1]{Caicedo2006}.
Atomlessness is a separate property and is not part of that definition.

We use the following preservation theorem of Solovay
\cite[Theorem~7]{Solovay1971}, in the formulation recorded in
\cite[Fact~1.27]{Caicedo2006}.

\begin{theorem}[Solovay]\label{model:solovay}
If $\kappa$ is real-valued measurable, then forcing with a measure algebra
preserves its real-valued measurability.  In particular, in the extension
there is a $\kappa$-additive probability measure on the full power set of
$\kappa$ in that extension, vanishing on singletons.
\end{theorem}

Here a measure algebra is the Boolean algebra of measurable sets modulo
null sets of a probability space, with zero omitted when it is used as a
forcing notion.  The special case of Theorem~\ref{model:solovay} used
below also has a direct proof in
\cite[Section~2, Lemmas~2.1--2.4]{FuchinoGreenbergShelah2006}.
That construction measures every subset of $\kappa$ in the extension,
is countably additive, and makes unions of fewer than $\kappa$ null sets
null. These properties give the stated additivity: in a disjoint family
of size less than $\kappa$, only countably many members have positive
measure, and the union of the remaining members is null.

\begin{proposition}\label{model:positive}
If ZFC with a measurable cardinal is consistent, then ZFC with a
real-valued measurable continuum is consistent.
\end{proposition}

\begin{proof}
Work in a ground model of ZFC with a measurable cardinal $\kappa$.
A $\kappa$-complete nonprincipal ultrafilter on $\kappa$ gives a
$\{0,1\}$-valued probability measure witnessing its real-valued
measurability.

We first record the cardinal arithmetic needed to control the continuum.
The cardinal $\kappa$ is regular: sets of size less than $\kappa$ are
null, so a cofinal union of fewer than $\kappa$ such sets cannot cover
$\kappa$.  It is also a strong limit.  Otherwise, for some
$\lambda<\kappa$ there would be an injection
$e:\kappa\to\mathcal P(\lambda)$.  For each $\xi<\lambda$, the ultrafilter
chooses one of the two sets on which the statement $\xi\in e(\alpha)$ is
constant.  Their intersection belongs to the ultrafilter by
$\kappa$-completeness, although injectivity makes that intersection have
at most one element.  This is impossible.  Every countable sequence of
ordinals below $\kappa$ is therefore bounded, and
\[
  \kappa^{\aleph_0}
   \leq \sum_{\lambda<\kappa}\lambda^{\aleph_0}
   \leq \kappa.
\]
Here $\lambda^{\aleph_0}\leq
2^{\max(\lambda,\aleph_0)}<\kappa$ for every $\lambda<\kappa$.
The reverse inequality is immediate, so
$\kappa^{\aleph_0}=\kappa$.

Let $\mathbb B_\kappa$ be the measure algebra of the fair-bit product on
$2^{\kappa\times\mathbb N}$.  We use the probability space generated by
finite-coordinate cylinders, or its completion; the quotient algebra is
the same.  Every element of this algebra has a representative determined
by countably many coordinates.  Such a representative is the inverse
image of a Borel set in a countable product.  Thus
\[
 |\mathbb B_\kappa|
 \leq \kappa^{\aleph_0}\,2^{\aleph_0}=\kappa
\]
in the ground model.  An antichain consists of disjoint nonzero elements
modulo null sets and hence is countable: for each positive integer $n$,
only finitely many of its elements can have measure greater than $1/n$.
Consequently the forcing has the countable chain condition and preserves
cardinals
\cite[Chapter~VII, Theorem~5.10]{Kunen1980}.

The generic coordinate sequences give $\kappa$ distinct reals.  Indeed,
for distinct $\alpha,\beta<\kappa$, the event that their coordinate
sequences agree has product measure zero. For the reverse bound, every real
has a name specified by countably many antichains deciding its bits.  Each antichain is countable, so at most
$|\mathbb B_\kappa|^{\aleph_0}\leq\kappa$ such deciding-antichain descriptions
suffice to name all reals
\cite[Chapter~VII, Definition~5.11 and Lemmas~5.12--5.13]{Kunen1980}.
Hence the extension satisfies $\mathfrak c=\kappa$.

Theorem~\ref{model:solovay} now gives the required measure on
$\mathcal P(\mathfrak c)$ in the extension.  Set forcing preserves ZFC,
including Choice~\cite[Chapter~VII, Theorem~4.2]{Kunen1980}.
The definable forcing relation and truth theorem
\cite[Chapter~VII, Theorems~3.5--3.6]{Kunen1980} give the asserted
relative-consistency implication in their syntactic, or finite-fragment,
interpretation~\cite[Chapter~VII, Sections~1 and~9]{Kunen1980}:
a finite contradiction proof in the extension theory would yield one in
the ground theory. This requires no transitive set model of all ZFC.
\end{proof}

\subsection*{Completion of the relative-consistency argument}
Let $\mathrm{SQ}_{\mathbb R}$ and $\mathrm{SQ}_{\mathbb C}$ denote the
separable quotient assertions over the indicated scalar fields.
Applying Theorem~\ref{pos:theorem} and Proposition~\ref{ev:positive-transfer} to
Proposition~\ref{model:positive} yields, for each
$\mathbb F\in\{\mathbb R,\mathbb C\}$,
\[
 \operatorname{Con}(\mathrm{ZFC}+\text{there is a measurable cardinal})
 \ \Longrightarrow\
 \operatorname{Con}(\mathrm{ZFC}+\mathrm{SQ}_{\mathbb F}).
\]
On the other side, G\"odel's constructible-universe theorem
\cite{Godel1938}, in the formulation of
\cite[Chapter~VI, Corollary~4.9]{Kunen1980}, gives
$\operatorname{Con}(\mathrm{ZFC})\Longrightarrow
\operatorname{Con}(\mathrm{ZFC}+\mathrm{CH})$.
Theorem~\ref{twist:counterexample}, together with Corollary~\ref{ev:fields},
therefore gives
\[
 \operatorname{Con}(\mathrm{ZFC})
 \ \Longrightarrow\
 \operatorname{Con}(\mathrm{ZFC}+\neg\mathrm{SQ}_{\mathbb F})
 \qquad(\mathbb F=\mathbb R,\mathbb C).
\]
Under the measurable-cardinal consistency hypothesis both sides are
consistent with ZFC.  Therefore ZFC proves neither
$\mathrm{SQ}_{\mathbb F}$ nor its negation, for either scalar field.

\subsection{Independence at density $\aleph_1$}

For the restricted assertion, the positive consistency bound can be
lowered using the earlier small-density theorem of Saxon and
S\'anchez Ruiz \cite{SaxonSanchezRuiz1996}: every infinite-dimensional
real Banach space of norm density less than $\mathfrak b$ has a separable
infinite-dimensional quotient; see \cite[Theorem~2.1]{Brech2025} for
its statement and a proof sketch. We explain the set-theoretic hypothesis that
makes this theorem apply at density $\aleph_1$.

The axiom $\mathrm{MA}_{\aleph_1}$ states that, in any nonempty partial order with
the countable chain condition, every collection of at most $\aleph_1$
dense sets is met by a filter. Full Martin's axiom, denoted by
$\mathrm{MA}$, requires this conclusion for every collection of fewer
than $\mathfrak c$ dense sets. The axiom $\mathrm{MA}_{\aleph_1}$ implies
$\mathfrak b>\aleph_1$.
To see this, fix a family $\mathcal A\subset\N^{\N}$ of cardinality at
most $\aleph_1$. A condition is a pair $(s,A)$, where $s$ is a finite
sequence of natural numbers and $A$ is a finite subset of $\mathcal A$.
A stronger condition $(t,B)$ extends $s$, contains $A$ in $B$, and
satisfies $t(n)\ge f(n)$ for every $f\in A$ and every new coordinate
$n$ of $t$. Conditions with a fixed stem $s$ have common extensions for
every finite subcollection: retain $s$ and take the union of their
finite sets. There are only countably many stems, so the partial order
is $\sigma$-centered and in particular has the countable chain condition.

For each $f\in\mathcal A$, the conditions with $f\in A$ form a dense
set. For each $n$, the conditions whose stem has length at least $n$
also form a dense set, since finitely many functions can be dominated
when a stem is extended. A filter meeting all these sets has a union
of stems $g\in\N^{\N}$. Once a condition has included $f$ in its finite
set, every common strengthening with a sufficiently long stem dominates
$f$ on the new coordinates. Directedness of the filter therefore gives
$f\leq^*g$. Thus $\mathcal A$ is bounded, proving
$\mathfrak b>\aleph_1$.

\begin{proof}[Proof of Corollary~\ref{intro:small-density}]
Starting with a ground model of ZFC and GCH, the forcing construction in
\cite[Chapter~VIII, Theorem~6.3]{Kunen1980} gives a model of Martin's
axiom with $\mathfrak c=\aleph_2$. G\"odel's constructible-universe theorem
and the relative-consistency interpretation of forcing therefore give
\[
 \operatorname{Con}(\mathrm{ZFC})\Longrightarrow
 \operatorname{Con}(\mathrm{ZFC}+\mathrm{MA}+\mathfrak c=\aleph_2).
\]
This model satisfies $\mathrm{MA}_{\aleph_1}$, hence
$\mathfrak b>\aleph_1$ by the preceding argument. The small-density
theorem gives a separable infinite-dimensional quotient for every real
Banach space of norm density $\aleph_1$. A complex Banach space has the
same norm density as its underlying real space, so
Proposition~\ref{ev:positive-transfer} gives the corresponding complex
conclusion.

On the other side, consistency of ZFC gives a CH model, and
Corollary~\ref{twist:density} supplies real and complex counterexamples
of norm density $\aleph_1$ there. Consequently, for each field
$\mathbb F\in\{\R,\C\}$,
\[
 \operatorname{Con}(\mathrm{ZFC})\Longrightarrow
 \begin{cases}
 \operatorname{Con}(\mathrm{ZFC}+\SQ^{\aleph_1}_{\mathbb F}),\\
 \operatorname{Con}(\mathrm{ZFC}+\neg\SQ^{\aleph_1}_{\mathbb F}).
 \end{cases}
\]
This proves the asserted independence with no large-cardinal consistency
hypothesis.
\end{proof}

